% !TeX root = ../../Skript.tex
\cohead{\Large\textbf{Potenzgesetze}}
\fakesubsection{Potenzgesetze}
Für Exponenten aus den natürlichen Zahlen \(n,\ m\in\N^*\) und Basen aus den reellen Zahlen \(a,\ b\in\R\) ergeben sich aus der Definition der Potenz die folgenden Gesetze:
\begin{tcolorbox}[title=Definition]\centering
    \(\textcolor{loestc}{a^n=\underbrace{a\cdot a\cdot a\cdot\dots\cdot a}_{n\text{-mal}}}\)
\end{tcolorbox}
Für gleiche Basis \(a\):

\begin{minipage}{\textwidth}
    \adjustbox{valign=t, padding=0ex 0ex 2ex 0ex}{\begin{minipage}{\textwidth/\real{3}-2ex}%
        \begin{tcolorbox}[height=1.5cm, valign=center, halign=center]
            \(a^n\cdot a^m=\ \)\textcolor{loestc}{\(a^{n+m}\)}
        \end{tcolorbox}
    \end{minipage}}%
    \adjustbox{valign=t, padding=1ex 0ex 1ex 0ex}{\begin{minipage}{\textwidth/\real{3}-2ex}
        \begin{tcolorbox}[height=1.5cm, valign=center, halign=center]
            \(\frac{a^n}{a^m}=\ \)\textcolor{loestc}{\(a^{n-m}\)}
        \end{tcolorbox}
    \end{minipage}}%
    \adjustbox{valign=t, padding=2ex 0ex 0ex 0ex}{\begin{minipage}{\textwidth/\real{3}-2ex}
        \begin{tcolorbox}[height=1.5cm, valign=center, halign=center]
            \((a^n)^m=\ \)\textcolor{loestc}{\(a^{n\cdot m}\)}
        \end{tcolorbox}
    \end{minipage}}%
\end{minipage}

\bigskip

Für gleiche Exponenten \(n\):

\begin{minipage}{\textwidth}
    \adjustbox{valign=t, padding=0ex 0ex 3ex 0ex}{\begin{minipage}{.5\textwidth-3ex}%
            \begin{tcolorbox}[height=1.5cm, valign=center, halign=center]
                \(a^n\cdot b^n=\ \)\textcolor{loestc}{\((a\cdot b)^n\)}
            \end{tcolorbox}
    \end{minipage}}%
    \adjustbox{valign=t, padding=3ex 0ex 0ex 0ex}{\begin{minipage}{.5\textwidth-3ex}
            \begin{tcolorbox}[height=1.5cm, valign=center, halign=center]
                \(\frac{a^n}{b^n}=\ \)\textcolor{loestc}{\(\left(\frac{a}{b}\right)^n\)}
            \end{tcolorbox}
    \end{minipage}}%
\end{minipage}